\section{Experiments}
\label{sec:experiments}

We evaluate MPI on the nine D4RL locomotion tasks formed by HalfCheetah,
Hopper, and Walker2d with medium, medium-replay, and medium-expert datasets
\citep{fu2020d4rl}. Base implementations follow CORL configurations
\citep{tarasov2022corl}; scores are D4RL normalized returns. Additional
per-task results and diagnostic definitions are in
Appendix~\ref{app:experiments}.

\paragraph{Fixed-budget protocol and estimand.}
We hold a total nominal budget $T$ fixed, split it across $K$ actor hops as
$h=T/K$, and sweep
\begin{equation*}
\begin{split}
T\in\{&0.05,0.1,0.2,0.4,0.7,1.5,2.5,\\
      &4,7,10,12,14,17,20\}.
\end{split}
\end{equation*}
over all nine tasks with two paired seeds. The sweep is complete over this
specified grid and contains $9\times14\times2$ runs per method. It is designed
as a high-resolution stability phase diagram rather than a final benchmark
alone. For $K=1$, Proposition~\ref{prop:spi_equiv} maps $T$ to the normalized
TD3+BC coefficient $\alpha=2T$.

We compare TD3+BC ($K=1$), proximal MPI with $K\in\{2,3\}$ (MPI-Prox), and a
scale-matched projected linearized target with $K\in\{2,3\}$ (MPI-Lin):
\begin{equation}
 a_{k+1}^{\rm tar}=\Pi_{\mathcal A}\!\left[
 a_k+d\frac{T}{K}\frac{\nabla_a\widehat Q(s,a_k)}{C_k}\right].
 \label{eq:linearized_target}
\end{equation}
Here $a_0$ is the sampled dataset action. At each hop, MPI-Lin detaches the
current reference, recomputes the $\widehat Q$ gradient and $C_k$, clips the
target to the action box, and takes one mean-squared-error optimizer step with
the corresponding persistent actor. It uses the same first-hop target and
final-hop evaluation routing as MPI-Prox. The factor $d$ matches the local
action-space coefficient induced by the action-averaged metric in
Eq.~\eqref{eq:mean_metric}. All conditions match datasets, critic updates, and
the horizon of $10^6$ critic updates; checkpoint diagnostics use the final
checkpoint. A $K$-hop condition performs $K$ actor optimizer steps at each
delayed policy update, and $C_k$ is recomputed at each reference. The primary
estimand is the end-to-end effect of the fixed-budget refinement bundle, which
jointly changes local step size, re-centering frequency, actor compute, and
target/deployment dataflow.

The analyses test whether this bundle reduces large-budget low-return
outcomes, whether the effect persists under the linearized realization, how
first-hop target-policy displacement differs from final-policy reach, and which
critic-validity diagnostics characterize failed training trajectories.

\subsection{Fixed-budget splitting broadens the stability envelope}

Figure~\ref{fig:stability} and Table~\ref{tab:stability} summarize the fixed-budget sweep. We descriptively call a high-budget cell collapsed when its two-seed mean score is below $20$; this is a locomotion failure indicator, not a theoretical threshold. MPI-Prox with $K=3$ has the strongest high-budget average and reduces the seed-mean collapse count from $35/63$ to $17/63$. MPI-Lin with $K=3$ reduces the same count to $24/63$, compared with $29/63$ for two linearized hops, and raises the high-budget average from $36.68$ to $43.91$. The raw-seed counts confirm the lower failure rate of both three-hop realizations. The proximal $K=2$ to $K=3$ mean improvement for $T\ge4$ is $+10.76$, whereas the median is only $+0.50$; the corresponding linearized differences are $+7.23$ and $+0.86$. Thus additional hops primarily rescue low-return outcomes rather than uniformly shifting settings. Appendix~\ref{app:results} reports sensitivity to thresholds $0,10,20,30,$ and $40$.

\begin{figure*}[t]
\centering
\includegraphics[width=\textwidth]{figures/stability_envelope.pdf}
\caption{Fixed-budget sweep over nine locomotion tasks. (a) Mean score over
tasks. (b) Collapse rates for $T\ge4$, reported both by environment--budget
cell and by seeded run. Fixed-budget substepping broadens the stable
high-budget region for both proximal and scale-matched linearized realizations.}
\label{fig:stability}
\end{figure*}

\begin{table}[t]
\centering
\caption{Aggregate fixed-budget results. High denotes $T\ge4$ and tail denotes
$T\ge10$. Collapse entries are seed-mean cells / raw seeded runs.}
\label{tab:stability}
\footnotesize
\begin{tabular}{lrrrr}
\toprule
Method & $T\le2.5$ & High & Tail & Collapse\\
\midrule
TD3+BC & 68.35 & 25.97 & 20.46 & 35 / 68\\
MPI-Prox, $K=2$ & \textbf{71.62} & 42.02 & 31.49 & 25 / 53\\
MPI-Prox, $K=3$ & 70.56 & \textbf{52.78} & \textbf{43.55} & \textbf{17 / 41}\\
MPI-Lin, $K=2$ & 69.02 & 36.68 & 27.44 & 29 / 58\\
MPI-Lin, $K=3$ & 68.54 & 43.91 & 33.12 & 24 / 53\\
\bottomrule
\end{tabular}
\end{table}

\paragraph{Task-dependent budget response.}
Environment-level curves exhibit recurring response motifs rather than a
scalar task-specific optimum. HalfCheetah medium and replay are comparatively
smooth, Hopper-replay has a wide high-return plateau, Hopper-medium and
Walker2d-medium show sharp, nonmonotone transitions, and expert
datasets show asymmetric high-budget failure. These motifs are descriptive,
not disjoint task classes.
In the matched Prox-2 audit, the current-state critic-coupled displacement
proxy organizes many of these transitions more directly than nominal $T$
(Appendix Figure~\ref{fig:environment_sensitivity}). HalfCheetah-medium remains
stable as $\widehat D_{\rm critic}$ grows from $.033$ to $.056$, whereas
Hopper-expert falls from $101.9$ to $1.8$ as the proxy grows from $.052$ to
$.907$. At comparable proxy values near $.1$, however, returns still differ
substantially across tasks. Nominal budget is therefore an indirect control
mapped through a task- and trajectory-dependent actor path and critic
evolution. Reversals and cross-task dispersion prevent these curves from being
summarized by a universal critic-valid radius; we use realized movement only
to organize within-task transitions.

\paragraph{Targeted $K=4$ extension.}
A separate-machine, four-seed MPI-Prox run on Hopper-medium shifts the high-return window to larger budgets. At $T=17$, the $K=3$ two-seed mean is $0.7$, whereas $K=4$ reaches $100.9\pm1.7$. The $K=4$ means at $T=\{10,12,14,17\}$ are $100.4,101.2,101.1,$ and $100.9$, with standard deviations at most $1.7$. At $T=20$, however, the four scores $99.7,51.0,0.7,$ and $15.8$ yield $41.8\pm44.0$; at $T=7$, $K=4$ scores $84.8\pm8.0$ versus $94.4$ for $K=3$. These per-seed results support a $K$-dependent, nonmonotone stability window with a sharp variance increase near $T/K=5$, not a verified $T_{\rm frontier}\propto K$ law. Changing $K$ jointly changes $T/K$, re-centering, compute, normalization, and routing, so this result does not identify a local-step mechanism (Appendix Table~\ref{tab:hopper_k4}). The per-seed scores have been transferred, but the external raw evaluation files, configurations, and checkpoints have not; we therefore exclude this targeted run from the complete-grid aggregate.

The completed linearized sweep shows that the stability gain survives linear-target regression and is therefore not unique to the proximal actor objective. Three linearized hops exceed TD3+BC by $17.94$ high-budget points and reduce collapse cells from $35$ to $24$; relative to two linearized hops, the average gain is $7.23$ points and the collapse count falls from $29$ to $24$. MPI-Prox $K=3$ remains stronger: over 63 paired high-budget cells, MPI-Lin $K=3$ wins by more than one point in 8, 21 are within one point, and MPI-Prox wins in 34; the mean and median Lin-minus-Prox differences are $-8.87$ and $-1.69$. Because the actor objectives differ, these remain comparisons of complete training procedures rather than isolated integrator effects.

\subsection{Large-budget behavior is nonlocal and critic dependent}

At $T=20$, the raw linearized target reaches the action box in
$94$--$95\%$ of states, and the effective action-step coefficient
$h_{\rm eff}$ is only $14\%$ and $9\%$ of the nominal $dT/K$ coefficient in
the two hops. These are medians over 18 environment--seed runs per hop. The
practical operation is projected target generation followed by amortized neural
regression under a moving critic. The sweep measures the resulting end-to-end
stability envelope; local truncation error lies outside this estimand.

Figure~\ref{fig:failure} shows the common failure signature. Environment
return supplies the external label: collapsed runs score below $20$. These
runs have orders-of-magnitude larger TD-error tails, larger critic magnitudes,
and larger realized movements. The controlled common critic is the frozen
final TD3+BC $T=1$ critic matched by environment and seed. A run-total gain
sums the mean $Q$ change across actor hops on the shared state batch. For
collapsed MPI-Lin runs, its median under the jointly trained critic is
$3.33\times10^{11}$; under the common critic it is $-1.42$, compared with
$+2.11$ for stable runs. For MPI-Prox, the common-critic gain changes from
$+2.51$ in stable runs to $-1.06$ in collapsed runs. This sign reversal
diagnoses cross-critic disagreement consistent with
Eq.~\eqref{eq:critic_error_margin}; environment return remains the criterion
for policy failure.

An exploratory CPU audit supplies a simulator reference for critic calibration.
Across 60 MPI-Prox checkpoints, the median signed error
$\widehat Q_1(s,a_K)-Q_{\rm MC}^{\rho_1}(s,a_K)$ is
$+1.70\times10^{12}$ in five collapsed cells and $-38.7$ in 55 stable cells,
corroborating severe overestimation in the collapsed screening cells. The
action-ranking subset is sparse: only 69 of 720 cloned states have both action
differences outside a one-return-unit dead zone, and 29 of those signs
disagree. We therefore use this audit as checkpoint-level calibration evidence
rather than as a routing or action-ranking estimand.

\begin{figure*}[t]
\centering
\includegraphics[width=\textwidth]{figures/failure_signature.pdf}
\caption{Median diagnostics over the matched
$T\in\{4,7,10,14,20\}$ audit. Proximal and linearized realizations share a
diagnostic pattern consistent with critic miscalibration: critic errors and
realized movements increase sharply, and run-total gains under a common critic become
negative. Environment return labels failure; the common critic diagnoses
cross-critic disagreement.}
\label{fig:failure}
\end{figure*}

\subsection{Target-policy displacement and final-policy reach}

At the same nominal budget, MPI-Prox realizes a smaller current-state first-hop
displacement proxy than MPI-Lin once $T\ge7$ (Appendix Figure~\ref{fig:exposure}a). At
$T=20$, their median $\widehat D_{\rm critic}$ values are $0.195$ and $0.323$.
Across 90 paired seeded comparisons at $T\in\{4,7,10,14,20\}$, exactly one
realization collapses in 11 cases. In all 11, the stable realization has both
smaller $\widehat D_{\rm critic}$ and smaller $\widehat D_{\rm final}$.
These paired outcomes associate smaller realized movement with the stable
optimization trajectory. Because the proxy is measured on current validation
states rather than the Polyak target policy on next states, this result
organizes the observed failure transition without identifying a causal route
through the TD target.

We directly measure Eq.~\eqref{eq:direct_target_displacement} on all 270
audited cells. The current-state proxy closely tracks deterministic
target-policy displacement ($r=.9999$). MPI-Prox $K=3$ has lower target-policy
displacement than TD3+BC in 88 of 90 paired runs (median ratio $.658$), while
the current-state $\widehat D_{\rm final}$ proxy is lower in only 36 of 90
(median ratio $1.057$). Thus the three-hop refinement can reduce movement in
the target-policy branch without merely shrinking its measured final-policy
reach. Against $K=2$, $K=3$ lowers target-policy
displacement in 84 of 90 runs but loses by more than one score point in 17 of
those 84, showing that target-policy displacement alone does not determine
performance (Appendix Table~\ref{tab:direct_exposure}).

A targeted $K=3$ study on Hopper-medium reveals two modes of decoupling
(Appendix Figure~\ref{fig:exposure}b). For $T=4$--$10$, $K=2$ and $K=3$ have nearly
identical $\widehat D_{\rm critic}$, whereas $K=3$ increases
$\widehat D_{\rm final}$ and score;
at $T=10$, the scores are $58.12$ and $85.35$. In the observed transition
region from stable to collapsed scores ($T=12,14$), $K=3$ instead contracts
both displacements and avoids the low-return outcome.
At $T=20$, both enter severe projection and saturation and collapse. The third
hop changes target-policy movement, re-centering frequency, and the
optimization trajectory rather than the formal order of the local update.

\subsection{Targeted component controls}

We next hold the actor trajectory fixed and change only the backup route. Two actor-detached shadow critics start from the same parameters and receive identical minibatches and target-smoothing noise; one bootstraps through the first-hop target actor and the other through the final-hop target actor. The study covers Hopper-medium at $T\in\{10,20\}$ and Walker2d-medium at $T\in\{4,7\}$, with two seeds and four checkpoints. Relative to a common simulator return under the effective first-route backup continuation, the final-route shadow has greater state-median symmetric relative error in all eight run-level averages. The final-minus-first route delta has mean $.124$, median $.013$, a run-bootstrap mean 95\% interval $[.008,.323]$, and an exact two-sided sign-flip $p=.0078$. Because the final-route critic targets a different continuation, this contrast measures route-induced alignment with the primary first-route value estimand, not each shadow critic's calibration to its own continuation. Neither shadow critic drives the actors, so the intervention does not estimate a routing effect on return.

Two additional eight-run controls probe simpler explanations (Appendix Table~\ref{tab:targeted_controls}). With three actor optimizer calls in every condition, sequential $K=3$, a fixed-reference chain, and a direct-final control have means $71.4,70.0,$ and $59.4$ and collapse counts $2/8,2/8,$ and $3/8$. Fixed-reference refinement therefore matches sequential re-centering in this targeted sample, so the experiment does not establish re-centering alone as the source of the sweep-wide gain. Slowing only the $K=1$ actor-target Polyak rate by a factor of three changes the matched collapse count from $6/8$ to $7/8$ and does not reproduce the $K=3$ rescue; actual displacement under this lag intervention was not audited. Together, the controls support first-hop backup routing while retaining the implemented bundle, rather than any isolated component, as the main estimand.

\subsection{Discussion and limitations}

Both three-hop realizations reduce large-budget failures, although MPI-Prox remains stronger overall. Direct measurements show that lower first-hop target movement need not shrink final-policy reach, and fixed-trajectory final-route shadows are less aligned with the common first-route value estimand in all eight targeted runs. Fixed-reference refinement matches sequential re-centering, while slower target averaging does not rescue the selected one-hop cells. Thus the primary estimate remains the implemented bundle: $K$ changes local step, re-centering, compute, normalization, and routing together. The fourth-hop study likewise relocates rather than uniformly improves the high-return window.

The complete grid covers TD3+BC on nine locomotion datasets with two paired seeds per cell. The separate-machine four-seed $K=4$ study lacks transferred raw artifacts. Collapse is descriptive, and the common-critic, simulator, and eight-run controls are targeted rather than exhaustive. The finite-step margin concerns the learned critic; the Wasserstein-flow account assumes small realized steps under a fixed critic.
